\chapter{交流实习}

\section{海外交流}
\subsection{总述——为什么要在海外实习}
在开始一段海外实习项目前，同学们应当明确自己的目的：单纯为了出国玩还是另有所求？如果单纯只为一段海外交流经历，为简历上添一栏内容，那么不论去哪里实习都可以，但是如果是为了未来境外升学、境内保研，那么一段过程顺利、收获颇多、对未来帮助极大的实习经验才是有意义的。

因此，这一章主要关注以申请PhD项目中获取一封较好的推荐信为目的的海外交流/实习经历，笔者将按照时间顺序一一讲述。

\subsection{实习机会获取}
首先明确一点，对于中国科学技术大学的本科生，进行海外交流需要在\href{vista.ustc.edu.cn}{全球化学习平台}网站上进行申报，申报获批后方可成行，同时这些项目都算作公派出境项目，包括校内派出和自主联系项目。同时由于算作公派出境

在全球化学习平台上可以查找到国际GVP和港澳台GVP项目，在其中可以报名参加，一般是中科大与其他学校的协议项目或是其他校内有所支持的项目，包括学期交换、学期访问、暑期交流等项目。对生科院学生而言，值得注意的有斯坦福大学暑期实习UGVR项目，香港科技大学暑期实习项目等，以往还有瑞士洛桑大学生命科学学院暑期实习项目，整体来说项目数量较少，值得前往的项目竞争也比较大。按照其他学校的学生路径的推荐，通过学期交换到美国其他学校进行一学期到一年的交换也可以在实验室进行实习，不过还未曾听说过科大生科院有这样做的，因为课程设置等原因，往往只有大四一学年可以留作毕业设计，但是申请季却是在大四上学期，确有不方便的地方。

\subsection{行前准备}
\subsection{工作生活}

\section{国内其他单位交流实习}